\cleardoublepage
\section{Place Value}
In our number system, every digit has a purpose. A $4$ does not always mean $4$. 
Sometimes it means $4000$. That depends on where it sits in the number. This 
idea is called place value. Each place in a number is ten times the value of 
the one to its right. A number like 4275 is not just four digits. It is a 
structure: $4000 + 200 + 70 + 5$. We do not read numbers one digit at a time. 
We read them as a system, where each part works because of its position.

Place value\index{place value} turns a few digits into large or small numbers. Take $9836$. 
The $8$ is not just an $8$. It is $800$, because it sits in hundreds of places. In 
5104, the 0 adds nothing but still matters. It shows that there are no tens. 
In $12345$, the $2$ sits in the thousands place and means $2000$. These examples may 
seem simple, but they show a key truth: digits mean more when you know 
where they stand.

Place value is more than vocabulary. It affects how we add, subtract, estimate, 
and think about numbers. When we know what each digit means, we can break 
numbers apart and put them back together. We stop guessing and start 
understanding. That shift helps us see math not as a list of rules but 
as a system that makes sense.

\textbf{Place value} tells us the value of each digit in a number, 
depending on its position. In our decimal system, each place is ten 
times the value of the place to its right.




\subsection*{Example: Understanding Place Value}

To truly grasp the concept of place value, let's take a closer look at the 
number $4275$. Here, the digit $4$ is found in the thousands place, which 
means it actually represents $4000$, not just $4$. Moving to the right, the 
$2$ is in the hundreds place, so it stands for $200$. The $7$, sitting in 
the tens place, is worth $70$, and finally, the $5$ in the ones place simply 
means $5$. This system, where each digit's value increases by a factor 
of ten as you move left, is what makes our decimal system so 
powerful and flexible.

To help visualize this, consider the following place value chart:

\begin{center}
\begin{tabular}{|c|c|}
\hline
\textbf{Place} & \textbf{Value} \\
\hline
Thousands & $1000$ \\
\hline
Hundreds & $100$ \\
\hline
Tens & $10$ \\
\hline
Ones & $1$ \\
\hline
\end{tabular}
\end{center}

Let's see how this works in practice. If we write $4275$ in expanded form, we get:
\[
4275 = 4000 + 200 + 70 + 5
\]

Now, imagine you are looking at the number $9836$. You might wonder, what is the 
value of the digit $8$ in this number? Since the $8$ is in the hundreds place, its 
value is $800$. This simple observation shows how the position of a digit 
determines its true value. Similarly, if you break down $5104$ into its 
expanded form, you see it as $5000 + 100 + 0 + 4$, making it clear what each 
digit stands for. In another example, the digit $2$ in $12345$ is in the 
thousands place, so it represents $2000$. Understanding these positions 
helps us see the true value of every digit in a number, which is a fundamental 
skill in mathematics.

\newpage
\subsection{Hand-drawing 1}
This is where you will see a hand-drawn answer to the previous question in \textbf{English}.
\begin{figure}[htbp]
  \centering
  \includegraphics[width=\linewidth]{example-image} % TODO: Update caption for actual content
  \caption{A hand-drawn answer in English.}
\end{figure}

\newpage
\subsection{Hand-drawing 2}
This is where you will see a hand-drawn answer to the previous question in \textbf{Korean}.
\begin{figure}[htbp]
  \centering
  \includegraphics[width=\linewidth]{example-image} % TODO: Update caption for actual content
  \caption{A hand-drawn answer in Korean.}
\end{figure}